\documentclass[10pt,a4paper]{amsart}
\usepackage[margin=19mm]{geometry}
\usepackage{amsmath,amssymb,mathtools}
\usepackage[scr=rsfs,cal=euler]{mathalfa}
\usepackage{xfrac}
\usepackage[unicode,hidelinks,bookmarks=false]{hyperref}
\newcommand{\ie}{i.e.\ }
\newcommand{\cf}{cf.\ }
\newcommand{\resp}{respectively\ }
\newcommand{\Subsection}{\subsection}
\providecommand{\nobreakdash}{\nobreak\hbox{-}}
\setlength{\emergencystretch}{2em}
\multlinegap0pt
\usepackage{amssymb}
\newtheorem{theorem}{Theorem}
\newtheorem{lemma}{Lemma}
\title{Avoidance Criteria for Normal Holomorphic Curves on Complex Projective Space}
\author{Gopal Datt, Rahul Gogoi \MakeLowercase {and} Kushal Lalwani}
\date{Source: arXiv:2605.07583v1 (CC BY 4.0)}
\begin{document}
\maketitle

\noindent\emph{Source and licence.} Avoidance Criteria for Normal Holomorphic Curves on Complex Projective Space. Version arXiv:2605.07583v1 (CC BY 4.0), distributed by arXiv under the Creative Commons Attribution 4.0 International licence, http://creativecommons.org/licenses/by/4.0/, as recorded in the arXiv licence metadata for this exact version and shown on the abstract page https://arxiv.org/abs/2605.07583v1. Full licence notice: \texttt{licenses/arxiv-2605.07583-CC-BY-4.0.txt}.\par\smallskip

\noindent\textit{Editorial scope.} Counted unit: the article's theorem T-ACFFHPS (source0.tex lines 167--174). The article's complete original proof of it is delivered with it (source0.tex lines 384--474, one \texttt{proof} environment of 91 lines, which the article places away from the statement). No mathematics is written, repaired or re-derived: the statement and the proof are the article's own lines, in the article's own notation, and no retained line is substituted or re-flowed.  Retained uncounted as the source-local context the proof needs: lines 268--275; lines 278--287; lines 292--301.  Nothing was omitted from any retained span, so the package carries no omission record.  The retained lines cite 4 works, whose entries are transcribed verbatim from the article's own bibliography source in the e-print and delivered with short editorial labels recorded in the item's reference mapping.  No retained line contains a figure environment, an image-inclusion command or drawing code, so no image is delivered and the package declares no asset dependency.  Editorial formatting only, in addition: an AMS \texttt{amsart} preamble that restates the article's own declarations and macros used by the retained lines, namely the environments \texttt{theorem}, \texttt{lemma} on separate counters, together with the standard packages those lines need.  The retained environments print in this document as Theorem 1, Lemma 1 in order of appearance, whereas the article prints the same items with the numbering of its own full document.  The complete native source of the article as distributed in the arXiv e-print for this exact version is bundled unchanged as \texttt{sources/200091/source0.tex}.
\begin{lemma}[{\cite[Theorem 2.5]{thai-fonmiscvahocs-2003}}] \label{L-Thai_Zalcman_Projective} 
Let $\mathcal{F}$ be a family of holomorphic mappings from a domain $\Omega\subset \mathbb{C}^m$ into $\mathbb{P}^n(\mathbb{C})$, where $m\in\mathbb N$. The family 
$\mathcal{F}$ is not normal on $\Omega$ if and only if there exist sequences $\{f_{\nu}\}_{\nu=1}^\infty \subset \mathcal{F}$, $\{z_{\nu}\}_{\nu=1}^\infty \subset \Omega$ with 
$z_{\nu} \to z_0 \in \Omega$ and $\{\rho_{\nu}\}$ with $\rho_{\nu}\searrow 0$ such that
$h_{\nu}(\zeta) := f_{\nu}(z_{\nu} + \rho_{\nu} \zeta)$
converges locally uniformly on $\mathbb{C}^m$ to a non-constant holomorphic mapping $h$ of $\mathbb{C}^m$ into 
$\mathbb{P}^n(\mathbb{C})$.
\end{lemma}

\begin{lemma}[{\cite[Theorem 1]{ere-apttfhc-1999}}] \label{L-Ere_omitting_hypersurface_general_position_implies_constant} 
Let $M \subseteq \mathbb{P}^n(\mathbb{C})$ be a closed subset  and say $\{Q_{j}\}_{j=1}^{2t+1}$ be a set of hypersurfaces in 
$\mathbb{P}^n(\mathbb{C})$ where
\begin{equation*}
M \cap \bigcap_{j \in I} Q_{j} = \varnothing
\quad \text{for every } I \subset \{1,2,\ldots,2t+1\},\ |I| = t+1.
\end{equation*}
Then any holomorphic mapping $f$ from $\mathbb{C}$ into
$M \setminus \bigcup_{j=1}^{2t+1} Q_{j}$ is constant.
\end{lemma}

\begin{lemma}[{\cite[Lemma 3.8]{liu-ofommitcps-2016}}] \label{L-Liu_hyperplanes_subgeneral_position_some_omit_some_intersect_finitely_implies_constant} 
Let $\{H_{j}\}_{j=1}^{2t+1}$ be hyperplanes in $\mathbb{P}^n(\mathbb{C})$ located in $t$\,-\,subgeneral 
position. Let $\hat{i}$ be an integer with $t+1 \leq \hat{i} \leq 2t+1$ and let $f : \mathbb{C} \to \mathbb{P}^n(\mathbb{C})$ be a 
holomorphic map. Assume that
\begin{enumerate}
\item given any $j=1,2,\dots,\hat{i}$, either $f(\mathbb{C}) \subset H_j$ or $f$ omits $H_j$ and
\item given any $j=\hat{i}+1,\dots,2t+1$, $\langle \mathbf f,H_j\rangle$ has atmost finitely many zeros in $\mathbb{C}$
\end{enumerate}
then $f$ is constant.
\end{lemma}

\begin{theorem}\label{T-ACFFHPS}
Let $\{H_i\}_{i=1}^{2n+1}$ be hyperplanes in general position in $\mathbb P^n(\mathbb C)$. Let $\mathcal F$ be a family of 
holomorphic curves from $\mathbb D$ to $\mathbb P^n(\mathbb C)$ such that $f^{-1}(H_i)=g^{-1}(H_i)$ on $\mathbb D$, for all 
$f,g\in\mathcal F$ and $i=1,2,\dots,2n+1$. Then there exists $M>0$ such that 
\begin{align*}\left(\frac{a-b\,|z|^2}{c+d\,|z|^2}\right)\|f'(z)\|_{FS}\leq M, \label{ACFFHPS} \tag{\ref*{T-ACFFHPS}} \end{align*}
for each $z\in \mathbb D$ and $f\in\mathcal F$, where $a,b,c,d\in\mathbb R$, $abc\neq 0$, $ab\geq 0$ and 
$c+d\,|z|^2\neq0$, for all $z\in\overline{\mathbb D}$.
\end{theorem}

\begin{proof}[{\bf Proof of Theorem~\ref{T-ACFFHPS}}]
Suppose that the statement fails to hold for $z\in \mathbb D$ and $f\in\mathcal F$. Then there exist $\{f_{\nu}\}_{\nu=1}^\infty\subset \mathcal F$, $\{z_{\nu}\}_{\nu=1}^\infty\subset\mathbb D$ such that 
\begin{align*}
    \left(\frac{a-b\,|z_{\nu}|^2}{c+d\,|z_{\nu}|^2}\right)\|f_{\nu}'(z)\|_{FS}\to\infty\ \text{ as }\ \nu\to\infty.\label{ACFFHPS.1} \tag{\ref*{T-ACFFHPS}.1}
\end{align*}
Thus \begin{align*}
    \left(\frac{\sqrt{a}-\sqrt{b}\,|z_{\nu}|}{c+d\,|z_{\nu}|^2}\right)\|f_{\nu}'(z_{\nu})\|_{FS}&=\left(\frac{a-b\,|z_{\nu}|^2}{(\sqrt{a}+\sqrt{b}\,|z_{\nu}|)(c+d\,|z_{\nu}|^2)}\right)\|f_{\nu}'(z_{\nu})\|_{FS}\\
    &\geq \frac{1}{\sqrt{a}+\sqrt{b}}\left(\frac{a-b|z_{\nu}|^2}{c+d|z_{\nu}|^2}\right)\|f_{\nu}'(z_{\nu})\|_{FS}.
\end{align*}
Hence \begin{align*}
\left(\frac{\sqrt{a}-\sqrt{b}\,|z_{\nu}|}{c+d\,|z_{\nu}|^2}\right)\|f_{\nu}'(z_{\nu})\|_{FS}\to\infty\ \text{ as }\ \nu\to\infty. \label{ACFFHPS.2} \tag{\ref*{T-ACFFHPS}.2} \end{align*}
Define $g_{\nu}(z):=f_{\nu}\left(z_{\nu}+\left(\frac{\sqrt{a}-\sqrt{b}\,|z_{\nu}|}{c+d\,|z_{\nu}|^2}\right)z\right)$, where $|z|<\dfrac{1-|z_{\nu}|}{\left(\frac{\sqrt{a}-\sqrt{b}\,|z_{\nu}|}{c+d\,|z_{\nu}|^2}\right)}$.
Then \begin{align*} \|g_{\nu}'(z)\|_{FS}&=\left(\frac{\sqrt{a}-\sqrt{b}\,|z_{\nu}|}{c+d\,|z_{\nu}|^2}\right)\left\|f'_{\nu}\left(z_{\nu}+\left(\frac{\sqrt{a}-\sqrt{b}\,|z_{\nu}|}{c+d\,|z_{\nu}|^2}\right)z\right)\right\|_{FS} 
\end{align*}
    In particular for $z=0$, we have
    \begin{align*}
    \|g_{\nu}'(0)\|_{FS}&= \left(\frac{\sqrt{a}-\sqrt{b}\,|z_{\nu}|}{c+d\,|z_{\nu}|^2}\right)\|f'_{\nu}(z_{\nu})\|_{FS} \to\infty\ \text{ as }\ \nu\to\infty.
\end{align*}
Now, from Marty's criterion \cite[p. 4]{ere-nhcfprtps-2007}, it follows that $\{g_{\nu}\}$ is not normal at origin. By Lemma 
\ref{L-Thai_Zalcman_Projective}, there is a subsequence of $\{g_{\nu}\}$ with $w_{\nu}\to 0$ and $\rho_{\nu}\searrow 0$
such that
\begin{align*}
    G_{\nu}(\zeta):=&\, g_{\nu}(w_{\nu}+\rho_{\nu}\zeta)\\
    =&\,f_{\nu}\left(z_{\nu}+\left(\frac{\sqrt{a}-\sqrt{b}\,|z_{\nu}|}{c+d\,|z_{\nu}|^2}\right)w_{\nu}+\left(\frac{\sqrt{a}-\sqrt{b}\,|z_{\nu}|}{c+d\,|z_{\nu}|^2}\right)\rho_{\nu}\zeta\right)
\end{align*}
where $\zeta\in\mathbb C$ satisfies $w_{\nu}+\rho_{\nu}\zeta\in\mathbb D$, converges locally 
uniformly on $\mathbb C$ to a non-constant holomorphic curve $G(\zeta)$ from $\mathbb C$ to $\mathbb P^n(\mathbb C)$. 
Considering a subsequence, we may now assume that $z_{\nu}\to z_0$, where $|z_0|\leq 1$. We define,
$$\widetilde{z_{\nu}}=z_{\nu}+\left(\frac{\sqrt{a}-\sqrt{b}\,|z_{\nu}|}{c+d\,|z_{\nu}|^2}\right)w_{\nu}\ \text{ and }\ \widetilde{\rho_{\nu}}=\left(\frac{\sqrt{a}-\sqrt{b}\,|z_{\nu}|}{c+d\,|z_{\nu}|^2}\right)\rho_{\nu}.$$ 
It is clear that $\widetilde{z_{\nu}}\to z_0$ and $\widetilde\rho_{\nu}\to 0$ as $\nu\to\infty$.\smallskip

Consider the hyperplanes $H_i$ for $i=1,2,\dots,2n+1$. Each $H_i$ is defined as,
\begin{align*}H_i=&\left\{[u_0:u_1:\cdots:u_n]\in\mathbb P^n(\mathbb C):\sum_{i=0}^n \alpha_I\,u^I=0\right\}\\ =&~\{[u_0:u_1:\cdots:u_n]\in\mathbb P^n(\mathbb C):L_i(\mathbf u)=0\}.\end{align*}

As given any $f,g\in \mathcal F$, $f$ and $g$ shares $H_i$, for all $i=1,2,\dots,2n+1$. Let 
$\{f_{\nu}\}_{\nu=1}^{\infty}\subset\mathcal F$. Then there exists $E_i\subset \mathbb D$, 
for all $i=1,2,\dots,2n+1$, $f_{\nu}^{-1}(H_i)=E_i$. \smallskip

If $E_i=\varnothing$, then $f_{\nu}^{-1}(H_i)=\varnothing$, for each $f_{\nu}\in\mathcal F$ and for all $i=1,2,\dots,2n+1$. If 
$z_0\in\mathbb D$, then for large $\nu$, $z_{\nu}\in D(z_0,r)$ for some $r>0$ with $D(z_0,r)\subset\mathbb D$. Now, 
$\widetilde{z_{\nu}}=z_{\nu}+\left(\frac{\sqrt{a}-\sqrt{b}\,|z_{\nu}|}{c+d\,|z_{\nu}|^2}\right)w_{\nu}$. As 
$w_{\nu}\to 0$ and $\left(\frac{\sqrt{a}-\sqrt{b}\,|z_{\nu}|}{c+d\,|z_{\nu}|^2}\right)$ stays bounded, we get 
$\widetilde{z_{\nu}}\to z_0\in\mathbb D$. Similarly $\widetilde \rho_{\nu}=\left(\frac{\sqrt{a}-\sqrt{b}\,|z_{\nu}|}{c+d\,|z_{\nu}|^2}\right)\rho_{\nu}\to 0$. 
Thus for each compact set $K\subset\mathbb C$, for large $\nu$, $\widetilde z_{\nu} +\widetilde \rho_{\nu} K\subset D(z_0,r)\subset \mathbb D$. \smallskip

As each $f_{\nu}(\mathbb D)$ omits all hyperplanes $H_i$, so $L_i(\mathbf{f}_{\nu}(z))\neq 0$, for all $z\in\mathbb D$. Hence
$L_i(\mathbf{G}_{\nu}(\zeta))\neq 0$, for all $\zeta\in\mathbb C$. Passing it to the limit and applying Hurwitz theorem,
\begin{itemize}
    \item $L_i(\mathbf{G}(\zeta))\equiv 0$ or,
    \item $L_i(\mathbf{G}(\zeta))\neq 0$ for all $\zeta\in\mathbb C$.
\end{itemize}


If $L_i(\mathbf{G}(\zeta))\equiv 0$ holds then there exists $i$ such that $G(\mathbb C)\subset H_i$ as 
the hyperplanes are in general position, their intersections with $H_i$ 
remains hyperplane in $H_i\equiv\mathbb P^{n-1}(\mathbb C)$. As there are $2n+1$ hyperplanes, there are atleast $2n$ hyperplanes in general position. Now observe 
$2n=2(n-1)+2$. So inductively, inside $H_i$, we still have atleast $2(n-1)+1$ hyperplanes. Thus $G$ must be constant. But by Lemma \ref{L-Thai_Zalcman_Projective}, 
$\|G'(0)\|_{FS}=1$, which is a contradiction. Hence $G(\mathbb C)\cap H_i=\varnothing$, for all $i=1,2,\dots,2n+1$.
Thus $G(\mathbb C)\subset \mathbb P^n(\mathbb C)\setminus \bigcup_{i=1}^{2n+1}H_i$. Also by Lemma \ref{L-Ere_omitting_hypersurface_general_position_implies_constant}, we get $G$ is constant, which is a contradiction. \smallskip

If $z_0\in\partial\mathbb D$, as $z_{\nu}\to z_0$, so $|z_{\nu}|\to 1$. Recall 
$\widetilde{z_{\nu}}=z_{\nu}+\left(\frac{\sqrt{a}-\sqrt{b}\,|z_{\nu}|}{c+d\,|z_{\nu}|^2}\right)w_{\nu}$ and 
$\widetilde \rho_{\nu}=\left(\frac{\sqrt{a}-\sqrt{b}\,|z_{\nu}|}{c+d\,|z_{\nu}|^2}\right)\rho_{\nu}$. As $w_{\nu}\to 0$, we 
still have $\widetilde{z_{\nu}}\to z_0\in\partial\mathbb D$. Thus for any compact set $K$ inside $\mathbb C$, 
$\widetilde{z_{\nu}}+\widetilde{\rho_{\nu}}K=z_{\nu}+\left(\frac{\sqrt{a}-\sqrt{b}\,|z_{\nu}|}{c+d\,|z_{\nu}|^2}\right)(w_{\nu}+\rho_{\nu} K)\subset\mathbb D$, 
the same procedure appears as above and $G$ turns out to 
constant by Lemma \ref{L-Ere_omitting_hypersurface_general_position_implies_constant}, which is again a contradiction. \smallskip

If $E_i\neq\varnothing$, then set $E_i=\{z\in\mathbb D:\langle\mathbf{f}_{\nu}(z),H_i\rangle=0\}$. Also here $E_i$'s
are discrete in $\mathbb D$, 
so given any compact set $K\subset \mathbb C$, the functions 
$\langle\mathbf{G}_{\nu}(\zeta),H_i\rangle$ have only finitely many zeros in $K$, for large $\nu$. \smallskip

Since $G_{\nu}\to G$ locally uniformly, we have $\langle\mathbf{G}_{\nu}(\zeta),H_i\rangle\to \langle\mathbf{G}(\zeta),H_i\rangle$ locally uniformly
on $\mathbb C$. Similarly, if $z_0\in\mathbb D$, then for each compact set $K\subset\mathbb C$, for large $\nu$,
$\widetilde z_{\nu} +\widetilde \rho_{\nu} K\subset D(z_0,r)\subset \mathbb D$.
Thus passing through a subsequence and applying Hurwitz's theorem, 
\begin{itemize}
    \item $G(\mathbb C)\subset H_i$ or,
    \item $\langle\mathbf{G}(\zeta),H_i\rangle$ has atmost finitely many zeros in $\mathbb C$.
\end{itemize}

Let $I=\{i:G(\mathbb C)\subset H_i\}$, $\widetilde{i}=|I|$. Since the hyperplanes are in general position in $\mathbb P^n(\mathbb C)$.
Thus atmost $n$ of them can have a common point. So $\widetilde{i}\leq n$. Define $J$ as the complement of the set $I$. Then $|J|=2n+1-\widetilde{i}\geq n+1$, for all 
$i\in J$, $\langle \mathbf{G}(\zeta),H_i\rangle$ has finitely many zeros. By Lemma 
\ref{L-Liu_hyperplanes_subgeneral_position_some_omit_some_intersect_finitely_implies_constant}, $G$ turns out to be a constant which is a contradiction. \smallskip

If $z_0\in\partial\mathbb D$, then again for any compact set $K$ inside $\mathbb C$, $$\widetilde{z_{\nu}}+\widetilde{\rho_{\nu}}K=z_{\nu}+\left(\frac{\sqrt{a}-\sqrt{b}\,|z_{\nu}|}{c+d\,|z_{\nu}|^2}\right)(w_{\nu}+\rho_{\nu} K)\subset\mathbb D$$ for large $\nu$ as above $G$ turns out to constant by Lemma
\ref{L-Liu_hyperplanes_subgeneral_position_some_omit_some_intersect_finitely_implies_constant}, 
which is again a contradiction. 
\end{proof}

\begin{thebibliography}{99}
\bibitem[ere-ap]{ere-apttfhc-1999} \label{ere-apttfhc-1999} Eremenko, A. (1999). \textit{A Picard type theorem for holomorphic curves}. Period. Math. Hung., 38, 39–42. \url{https://doi.org/10.1023/A:1004794914744}
\bibitem[ere-nh]{ere-nhcfprtps-2007} \label{ere-nhcfprtps-2007} Eremenko, A. (2007). \textit{Normal holomorphic curves from parabolic regions to projective spaces}. \url{https://doi.org/10.48550/arXiv.0710.1281}
\bibitem[liu-of]{liu-ofommitcps-2016} \label{liu-ofommitcps-2016} Liu, X., Pang, X., Yang, L. (2016). \textit{On families of meromorphic maps into the complex projective space}. Houston journal of mathematics. 42. 775-789. \url{MR3570710}
\bibitem[thai-f]{thai-fonmiscvahocs-2003} \label{thai-fonmiscvahocs-2003} Thai, D. D., Trang, P. N. T., Huong, P. D. (2003). \textit{Families of normal maps in several complex variables and hyperbolicity of complex spaces}. Complex Var. Theory Appl. 48:6 469–482. \url{MR 2004c:32003 Zbl 1036.32001}
\end{thebibliography}
\end{document}
